\documentclass[12pt]{article}

\usepackage[utf8]{inputenc}
\usepackage[T1]{fontenc}

\title{El lenguaje científico de la Biología}
\author{Kihara G. Rendon, Sara S. Pay, Diego A. Castro, Alessandro I. Molina}
\date{11/09/2026}

\begin{document}

\maketitle
\section{Autores del artículo de referencia}
\textbf{Artículo de referencia:} Bazán Delgado Anieska. (2017). \textit{LA BIOLOGÍA Y SU LENGUAJE CIENTÍFICO, UN ANÁLISIS LINGüÍSCO}. Revista Santiago.

Anieska Bazán Delgado es Doctora en Ciencias Pedagógicas y Profesora de Microbiología y Virología de la universidad de Oriente, Santiago de Cuba. Cuenta con una trayectoria destacada en eventos académicos y publicaciones en revistas. Su trabajo profesional combina la práctica docente en el área de la ciencias con la investigación educativa, centrandose en cómo las particularidades del discurso científico influyen en el aprendizaje y en la formación de los estudiantes.

\section{Introducción}
En la formación escolar y universitaria, el aprendizaje de las ciencias ha enfrentado barreras como la memorización pasiva o la falta de motivación. La autora aborda este problema centrándose en la complejidad del vocabulario científico. En su trabajo, integró los aspectos conceptuales y léxicos que diferenciaban este discurso del habla cotidiana del resto de disciplinas. Para quienes nos formamos en el área de la Biología, este enfoque resulta clave para comprender que el manejo preciso de la terminología no es un simple formalismo, sino es la herramienta fundamental para construir un pensamiento científico.

\section{Métodos}
Anieska revisó las funciones cognitivas y comunicativas del lenguaje para analizar la enseñanza cientifica. Explicó cómo se construye el léxico mediante palabras con raices griegas y latinas, abreviaciones y palabras prestadas de ramas hermanas a la Biología como la Química, la Física y la Geografía. La razón de usar este enfoqué didactico radicó en evitar que la enseñanza de la ciencia se redujera a memorizar lista de conceptos, enseñando en su lugar a interpretar códigos simbólicos en su contexto real.

\section{Resultados}
En un aspecto de forma léxica, la autora pudo señalar que los textos biológicos son caracterizados por emplear ciertos adjetivos con la función de sustantivoss y contrucciones pronominales que contribuyen a darle un caracter con un toque más impersonal. Un ejemplo frecuente de esto podria ser el uso frecuente del "se", siendo un recurso que permite expresar las ideas sin tener que hacer una referencia directa a una persona o objeto determindo. Asimismo, resaltó la importancia histórica de la nomenclatura propuesta por Linneo y de los epónimos, utilizados para reconocer y poder recordar los aportes cientificos como Golgi, Pasteur y Krebs, o para hacer referencia a personajes de la mitología.

En cuanto a la estructura sintáctica, los textos biológicos suelen presentar oraciones que permiten relacionar diferentes ideas y condensar una mayor cantidad de información. Además de recurrir a formas no personales del verbo, como puede ser el infinitivo, el gerundio y el participio, las cuales contribuyen a que la información se presente de manera más precisa, organizada y propia del lenguaje científico.

\section{Discusión}
La precisión en el uso de los términos y la objetividad son importantes para comunicar los resultados dentro de la comunidad académica. En biología, conocer el vocabulario especializado no solo permite describir un fenómeno, sino tambien comprender mejor lo que se está estudiando. Por esta razón, sería interesante incluir actividades de análisis etimológico en las prácticas de laboratorio de primer semestre. En vez de enfocarse solamente en memorizar los términos para los parciales, se podria trabajar con el origen de algunas raíces griegas y latinas para observar si esto ayuda a los estudiantes a comprender los conceptos y a mejorar su forma de redactar. De esta manera, aprender la estructura de estos términos podría ser una herramienta más para facilitar la comprension de los temas sin convertir el aprendizaje en una simple memorización.

\section{Conclusión}
La lingüística de la biología revela una dimensión del conocimiento científico que con frecuencia queda oculta detrás de los contenidos. La biología no se aprende únicamente mediante la observación de organismos, la experimentación o el estudio de modelos; también se aprende al dominar un sistema de palabras, estructuras, símbolos, convenciones y formas de argumetar. Ese sistema permite nombrar entidades, clasificar formas de vida, representar procesos, formular relaciones causales, comparar evidencias, y construir explicaciones. Por ello, el lenguaje no es un envoltorio externo de la ciencia biológica: es uno de los medios a través de los cuales la ciencia adquiere forma y puede ser compartida.

El análisis del lenguaje muestra que la terminologìa biológica contiene patrones morfológicos, raíces grecolatinas, derivaciones, abreviaciones y préstamos interdisciplinarios que pueden convertirse en recursos pedagógicos. La evidencia de Brown (2014) y Keen et al. (2025) sugiere que enseñar a interpretar la estructura de las palabras puede aportar beneficios al aprendizaje, aunque no sustituye la enseñanza conceptual. Del mismo modo, los estudios de Rector et al. (2013) muestran que la ambigüedad léxica puede afectar las explicaciones de evolución, mientras que trabajos sobre sintaxis y nominalización evidencian que comprender la gramática del discurso científico es necesario para leer textos biológicos complejos (Hao & Humphrey, 2019; Halliday & Martin, 1993).
\end{document}
